\chapter{再论微积分的基础}

\section{实数理论}

微积分的理论基础是在微积分发展到相当程度以后才建立起来的。我们在第一章一开始就指出，微积分既然是变量的数学，而变量又是一个极困难的概念，要想明确地理解它而不出逻辑上的矛盾，必须要等到具体的数学结果积累到相当丰富，把问题的各个方面都暴露到相当程度才行。所以，尽管在希腊时代就已经认识到诸如连续性、运动、无穷小等等是多么困难的概念，甚至欧多克萨斯关于比例的理论、关于不可通约量的理论与现代的实数理论颇为相近；在欧几里得关于圆面积的穷竭性理论中也看到了如何把极限问题化为不等式问题。但我们究竟不能说这些古希腊的著作就是今天通用的微积分教本所介绍的理论的前身，因为一些最根本的概念一直到19世纪中晚期才弄明白。而在希腊时代，数学还与种种思辨性的议论分不开。不仅毕达哥拉斯是这样，芝诺和亚里士多德更是这样。可是，科学（数学也是其一部分）还是在沿着探索宇宙的规律的大道上阔步前进。到了19世纪，这些问题在当时的水平上弄得相当清楚了。数学，特别是微积分，在说明宇宙的规律取得何等伟大的成就，相信前五章的内容已经足够给读者一个印象。既然如此，又何必回到历史的陈迹中，再去追忆当年呢？问题在于，数学还在前进。今天已不再是达朗贝尔“前进吧，您会得到信心”的时代。其实，自从所谓微积分的基础得到建立时起，整个数学又陷入更深刻的困难之中。这就是集合论的建立。论时间是19世纪七八十年代，主要的代表人物是康托尔（Georg Cantor），而且他的起点正是黎曼、狄利克雷关于三角级数的工作，我们在第五章一开始就介绍了一点点有关情况。今天我们再回到微积分的基础是因为，弄明白了这个基础，就发现它含有许多“成分”——用布尔巴基的语言来说是许多数学构造或称为 architecture。这些构造甚至规定了下一步数学将怎样走法。因此，本章的目的是从介绍微积分的基础使读者看到这些构造是些什么。其实在前几章中我们已在许多具体问题上与它们打过交道了。所以，我们希望是从再论微积分的基础中帮助读者能与更新、更现代的数学打交道。

\noindent\textbf{1. 有理数系怎样扩大成实数系}\quad
为微积分建立基础的工作现在都说是其算术化的过程。其原因在于极限理论过去实际上是以几何直观为基础的。例如，我们说
\[
 \lim_{x\to a}f(x)=b
\]
时，心目中总会想到当 $x$ 连续地沿 $x$ 轴趋向 $a$ 时，$f(x)$ 沿另一条轴趋向 $b$ 点。可是如果在 $x$ 轴没有一个点代表 $a$，或在另一条轴上没有一个点可以说是 $b$，则这样的极限式应如何理解？这种情况希腊人早就发现了，$x$ 轴上有一个点（距原点为单位边长正方形的对角线长的那一点）不能用一个数（希腊人讲数实际上是有理数）来表示，因此 $x$ 趋向这样的点，就不能说成是 $x\to a$。总之，我们的极限概念实际上有一个默认的基础：直线上的点与实数之间有一个一一对应。从几何上看，这似乎是不言而喻的，但是从逻辑上讲是有待证明的。实际上微积分基础中与这个没有明说的默认有关的许多基本的命题，如柯西序列必有极限，单调有界序列必有极限，其证明都有赖于此。

因此问题就在于如果只承认有理数，则作为数轴的直线上有一些点找不到对应的数，所以，有理数是有“空隙”的。只有把这些空隙都填补起来，极限理论才有可靠的基础。我们不能采取这样的办法，即缺什么数就补什么数：例如，有理数中既然没有表示单位边长正方形的对角线长的数，就补进一个 $\sqrt2$，然后补上例如 $\sqrt[3]7,e,\pi$ 等等。这样做何时才知道“所有的”空隙都补上了（连什么叫“所有的空隙”都说不清楚）？因此我们需要一个系统的方法补上一些数，使得以后再也没有必要、没有可能再添补新数。怎样知道是“再也没有必要、没有可能”的呢？对此也必须作出数学上很明确的回答。

在进行这项工作前，我们要弄清我们现在所处的地位以及我们想要做到的事。

我们所处的地位是：我们已经有了有理数系 $\mathbb Q$，它有以下的性质：

\noindent I. $\mathbb Q$ 是一个域，即在其中可以进行加、乘及其逆运算减、除（但不允许以0为分母）等四则运算，而且加法与乘法均适合交换律与结合律，二者在一起则适合分配律。

\noindent II. $\mathbb Q$ 是一个有序域，即对任意 $x,y\in\mathbb Q$，以下的三种关系必有一种且只有一种成立：$x<y$，$x=y$，或 $x>y$。

这种次序关系有以下几个性质：
\begin{align*}
 (\mathrm{II}_1)&\quad \text{若 }x\le y,\ y\le x,\text{ 则 }x=y;\\
 (\mathrm{II}_2)&\quad \text{若 }x\le y,\ y\le z,\text{ 则 }x\le z;\\
 (\mathrm{II}_3)&\quad \text{若 }x_1\le x_2,\ y_1\le y_2,\text{ 则 }x_1+y_1\le x_2+y_2;\\
 (\mathrm{II}_4)&\quad \text{若 }x\le y\text{ 而 }z>0,\text{ 则 }xz\le yz.
\end{align*}

\noindent III. $\mathbb Q$ 是一个阿基米德域，即满足以下的阿基米德公理：任给 $x,y$ 而 $x>0,y>0$，则必存在一个正整数 $n$，使 $y\le nx$。

把 I—III 当作公理来看，这里面其实是有问题的，例如对自然数确有一组人们公认的公理——佩亚诺（G. Peano，1858—1932）公理，而它并不就是 I。所以 $\mathbb Q$ 中的算术运算的许多性质应该是可以证明的而不应作为公理。这样说来，I 和 II 是否公理可以深入讨论，但我们完全不去接触这类问题而是把 I 和 II 当作已经明确不再讨论的命题。但是 III 确实是一个公理。它的重要性在于可以由它肯定，$\mathbb Q$ 中不会有无穷小 $\alpha$（不妨设 $\alpha>0$，因为所谓无穷小就是指非0的，而绝对值又小于任何有理数的“数”，所以 $\alpha\ne0$，若 $\alpha<0$ 则考虑 $-\alpha$）。因为若有无穷小 $\alpha$，则取 $y=1,x=\alpha$，由公理 III 可知必存在正整数 $n$ 使 $1\le n\alpha$，因此 $\alpha\ge1/n$ 而不可能小于有理数 $1/n$。同样也不会存在一个有理数 $\beta$ 是正无穷大，因为若 $\beta>$ 任意正有理数 $x$，$1/\beta<1/x$，而 $1/x$ 确实可以表示任意正有理数，所以 $1/\beta$ 就是无穷小。这个问题我们在第一章就已讨论过。

有理数系的最大“缺点”就是 $\mathbb Q$ 中的柯西序列不一定在 $\mathbb Q$ 中有极限，例如 $1,1.4,1.414,\ldots$ 就不能以任意有理数为极限。这就是 $\mathbb Q$ 中尚有“空隙”。正式的术语则说 $\mathbb Q$ 缺少完备性。因此我们的任务是：要对 $\mathbb Q$ 中补充一些数，使 $\mathbb Q$ 扩大成为实数系 $\mathbb R$，而要求

\noindent\textbf{1.}\ 给出一个补充 $\mathbb Q$ 的有系统的方法，使 $\mathbb R$ 不再有空隙，即有完备性。就是说我们上面说的柯西序列没有极限这种空隙要补上。这个补空的过程就称为完备化。我们不仅要问，完备化是否可能，还要问完备化的方式是否在某种意义下是唯一的？这个问题很重要。因此如果有不同的方法作完备化，而且作出了不同的实数系，则应该有不同的极限理论，有不同的微积分学。一旦出现这种情况，后果将是灾难性的。

\noindent\textbf{2.}\ $\mathbb Q$ 再补上一些数——称为无理数，说是“一些数”实在把完备化的内容说得太简单了，因为补上的新数要比原有的有理数多得多：原来的有理数只有可数无穷多个，而无理数是不可数无穷多个。这样构成了实数系 $\mathbb R$。但是尽管补上的新数多得多，原有的 $\mathbb Q$ 却在 $\mathbb R$ 中稠密。正是这个稠密性保证了新数仍然适合上面对 $\mathbb Q$ 的要求 I 和 II。因此我们确实可以承认它们都是数。但是 III 却还是在 $\mathbb R$ 上成立。因此，无穷小量与无穷大量都不是现在我们所理解的数，而它们是某种变量。

\noindent\textbf{3.}\ 有了实数系 $\mathbb R$，而且得知它已经是完备的，我们就要用它来证明微积分的一些基本的命题，例如单调有界数列有极限存在。我们下面会讲到，在有了完备性以后，不但不会因柯西序列没有极限而发现原有的数系有需要填补的空隙，而且也不会例如有界单调数列没有极限而产生空隙，最终我们会看到，在一定的要求之下，也不会因为任何问题再出现空隙。就是说，我们对问题的解决是一劳永逸的。

现在我们就按这个计划来建立实数系。

第一步是介绍完备化过程。

我们的基本思想是：柯西序列就是实数，这本来是很直观的，例如，凡讲到 $\pi$，人们心目中就一定是指一个序列，
\[
 \{a_n\}=\{3,3.1,3.14,3.141,3.1415,3.14159,\ldots\},
\]
即 $\pi$ 的无穷小数展开式的前有限项。它显然是一个柯西序列，因为例如其中第 $m$ 项 $a_m$ 是 $\pi$ 的展开式到小数点后 $m-1$ 位，第 $n$ 项 $a_n$ 则是展开式小数点后前 $n-1$ 位。设 $m<n$，则 $|a_m-a_n|$ 小数点后前 $m-1$ 位全都是0，因而
\[
 |a_m-a_n|\le \frac1{10^{m-1}},
\]
所以当 $m,n$ 充分大时，$|a_m-a_n|<\varepsilon$，$\varepsilon>0$ 是任意指定的有理数。所以它是柯西序列。为什么要说 $\varepsilon$ 是有理数呢？因为现在只有有理数才有定义。所以现在有理数的柯西序列定义应是 $\{a_n\},a_n\in\mathbb Q$，对任意正有理数 $\varepsilon>0$，必可找到正整数 $N(\varepsilon)$，使当 $m,n>N(\varepsilon)$ 时 $|a_m-a_n|<\varepsilon$。但在以后定义了实数以后，柯西序列中的 $\varepsilon>0$ 就可以改说是“任意正的实数 $\varepsilon>0$”，不过读者会看出即令不作这样的更改也没有影响。所以，下面我们就不再强调这一点了。我们想要强调的是，$\pi$ 还可以表示为另一个柯西序列
\[
 \{b_n\}=\{4,3.2,3.15,3.142,3.1416,3.14160,\ldots\}.
\]
所以，说一个柯西序列定义一个实数，并不恰当。不过 $a_n-b_n\to0$（读者可能会问，不是现在还没有定义极限吗？为什么能使用极限的记号呢？其实说极限没有定义是指序列、极限只要有一个不是有理数时都没有定义。现在我们的 $a_n,b_n,\varepsilon,\ldots$ 都是有理数，这时甚至使用极限的 $\varepsilon$—$N$ 定义都没有困难）。因此若有两个柯西序列 $\{a_n\},\{b_n\}$，适合 $a_n-b_n\to0$，就说它们是等价的，记作
\begin{equation}
 \{a_n\}\sim\{b_n\}\quad(\text{即指 }a_n-b_n\to0).
 \tag{1}
\end{equation}

\noindent\textbf{定义1}\quad
有理数的柯西序列按(1)式的等价类就是一个实数。

有的读者可能读过通常的实数理论，那里是以戴德金分割作为实数之定义的，何以这里要标新立异？何况现在我们的作法反而麻烦一些？理由在于今后我们将仿照由有理数构造实数的方法来构造许多函数空间，特别是 $L^p$ 空间。那时戴德金分割就用不上了。

于是，这个等价类中每一个具体的柯西序列都是这个等价类的代表元。在一个等价类中取哪一个柯西序列为代表元都是一样的。

有理数也是实数，那么哪一些等价类代表有理数呢？其实，若 $r$ 是一个有理数，令 $a_n=r$，得
\[
 \{r,r,\ldots,r,\ldots\}\quad(\text{记作 }\{r\})
\]
也是一个柯西序列，我们不妨称为常驻序列。我们不妨把 $r$ 本身与 $r$ 构成的常驻序列集合一一对应。因此，若一个等价类含有一个常驻序列 $\{r\}$，我们就认为这个等价类代表有理数 $r$。这里要注意，若一个等价类中含有常驻序列，则它不能含两个不同的常驻序列 $\{r\}$ 与 $\{s\},r\ne s$。事实上若 $r-s\to0$ 则必有 $r=s$，而与 $r\ne s$ 矛盾，所以 $\{r\}$ 与 $\{s\}$ 不可能在同一个等价类中。于是我们把柯西序列等价类的集合分为两类：I 是含有某个常驻序列 $\{r\}$ 的等价类之集合，II 类之元是不含常驻序列的等价类。

\noindent\textbf{定义2}\quad
I 类之元称为有理数，$\{r\}$ 对应于 $r$；II 类之元称为无理数。两类合并即得实数系 $\mathbb R$。

本来现在应该就来说明 $\mathbb R$ 为什么就是完备的，但在这以前，我们先来讨论实数的运算和次序。

第二步，实数的运算和次序关系。

实数的运算是比较容易定义的，设有两个实数 $a$ 与 $b$，并取其代表元 $\{a_n\}$ 和 $\{b_n\}$，则例如 $a+b$ 即定义为 $\{a_n+b_n\}$ 之等价类；$ab=\{a_nb_n\}$ 之等价类。我们要注意，$\{a_n+b_n\},\{a_nb_n\}$ 仍是有理数的柯西序列，所以这些运算是有定义的。加法零元0是包含常驻序列 $\{0,0,\ldots,0,\ldots\}$ 的等价类；乘法单位元则是包含常驻序列 $\{1,1,\ldots,1,\ldots\}$ 的等价类。$1/a$ 是 $\{1/a_n\}$ 的等价类，不过表面上应要求 $a_n\ne0$，而不仅是 $\{a_n\}\not\sim\{0,0,\ldots,0,\ldots\}$。但实际上，如果 $a\ne0$，则 $\{a_n\}$ 中只会含有最多有限多个 $a_n=0$。否则，若有无限多个 $a_{n_k}=0,k=1,2,\ldots,n_k\to+\infty$，则对任一 $\varepsilon$，因为当 $m,n>N(\varepsilon)$ 时 $|a_m-a_n|<\varepsilon$，我们只要取 $n=n_k>N$ 则有 $|a_m|=|a_m-a_{n_k}|<\varepsilon$，当 $m>N(\varepsilon)$ 时成立，即是说 $a=\{a_n\}\sim0$，这与 $a\ne0$ 矛盾。既然在 $\{a_n\}$ 中只有有限多个为0，可以从 $\{a_n\}$“砍掉”前面有限多项，余下来仍然是一个柯西序列，仍然与原来的 $\{a_n\}$ 等价，从而仍是 $a$ 的代表元。但现在一切 $a_n\ne0$，所以非零的 $a$ 一定有一个代表元 $\{a_n\}$ 其一切项均不为0。这样，$1/a=\{1/a_n\}$ 仍为柯西序列就很清楚了。总之，我们构造的新的数系：有理数柯西序列之等价类的集合，亦即实数系 $\mathbb R$ 仍是一个域（代数意义下的域）。

关于 $\mathbb R$ 是一个有序域的证明麻烦一些。我们可以很直观地想到，可以用 $a-b$ 为一正实数来表示 $a>b$，但为此对什么是正实数要认真论证一下。而且我们可以更广泛一些，分别讨论什么是 $a=0,a>0$ 和 $a<0$，这里的 $a$ 表示某柯西序列 $\{a_n\}$ 的等价类，而 $a_n$ 则是有理数。$a=0$ 就是适合 $\lim a_n=0$ 的序列，这里我们又借用了有理数范围内的极限语言，而如前面所说这是许可的。如果不如此我们就要说：对任意 $\varepsilon>0$，必有 $N(\varepsilon)$ 存在而当 $n>N(\varepsilon)$ 时，$|a_n|<\varepsilon$。$a>0$ 问题复杂一些，我们决不能说 $a>0$ 就是 $\{a_n\}$ 中一切 $a_n>0$。因为例如 $\{1/n\}\to0$ 但 $1/n>0$。这里我们要利用一个重要的事实，一个有理数 $s>0$ 时一定可以找到另一个 $r$ 使 $s>r>0$。可否把这个事实用于 $\{a_n\}$？我们说：$a=\{a_n\}>0$ 就是存在一个正有理数 $r=\{r\}$ 以及一个正整数 $N$，使当 $n>N$ 时，$a_n>r>0$。同样 $a=\{a_n\}<0$ 就是存在一个负有理数 $s=\{s\}$，以及一个正整数 $M$，使当 $n>M$ 时，$a_n<s<0$。这样的定义行不行，就看是否每一个实数 $a$（设用有理数的柯西序列 $\{a_n\}$ 来表示）都能无疑义地规定其符号。这是肯定的，因为 $\{a_n\}$ 是柯西序列，故对每个 $\varepsilon>0$ 必有 $N(\varepsilon)$ 存在，使当 $m,n\ge N(\varepsilon)$ 时，均有 $|a_m-a_n|<\varepsilon$。所以对一切 $m>N(\varepsilon)$，
\[
 a_m=a_{N}+(a_m-a_N),
\]
\[
 a_N-|a_m-a_N|\le a_m\le a_N+|a_m-a_N|.
\]
这样 $a_m$ 之符号将与 $a_N$ 的大小有密切关系。这样就出现两个可能性。其一是对于某一个 $\varepsilon>0$，$|a_N|>\varepsilon$。如果 $a_N$ 为正，则 $a_N>\varepsilon$，由左边的不等式看 $a_m>a_N-\varepsilon=r>0$，这样的 $\{a_n\}$ 符合前述 $a>0$ 之定义。如果 $a_N$ 为负，则 $a_N<-\varepsilon$，由右边的不等式看 $a_m<a_N+\varepsilon=s<0$，这样的 $\{a_n\}$ 符合前述 $a<0$ 的定义。第二个可能性是对一切 $\varepsilon$，按上述求出的 $a_N$ 均适合 $|a_N|\le\varepsilon$。这时当 $n>N(\varepsilon)$ 时，
\[
 |a_n|\le|a_N|+|a_n-a_N|<2\varepsilon.
\]
所以符合前述 $a=0$ 的定义。总之任一实数 $a$ 在 $a>0,a=0,a<0$ 三种情况中必居其一且仅居其一。实数 $a$ 既已可定义其符号，当然也就可以定义 $|a|$。关于 $\mathbb Q$ 为有序域的四条性质 $(\mathrm{II}_1)$—$(\mathrm{II}_4)$ 可以证明，对 $\mathbb R$ 也都适用。但是证明却可能比较冗长。

阿基米德公理对于 $\mathbb R$ 仍然是成立的。这一点也不来证明了。它的重要性下面就会看到。

第三步是关于实数系 $\mathbb R$ 的基本性质的证明。现在我们已经有了实数系了。此时不妨回顾一下，什么是实数？实数就是一个可以对之进行四则运算、并可比较大小，且适合阿基米德公理的对象。其实，我们对数的领悟已经有了一个飞跃。比方有理数，朴素的理解就是一个屈指可数的数。它的大小，一则可以用扳手指头的方法来比较（用抽象的数学语言来说，就是与“手指头的集合”一一对应），也可以通过几何方法用“量”来比较：圆内接多边形因为只是圆的一部分，部分小于整体，所以这些内接多边形面积小于圆面积（《几何原理》上就是这样讲的）。但是到了实数就一般地没有这种朴素的理解了。其四则运算化成了有理数柯西序列的运算，其大小则归结为与0的比较，即正、负号。读者们会问，实数诚然可以比较大小，但“它自己”有没有“大小”？上面说的思想上的飞跃在此：能比较大小就是有大小。所以圆面积不必看成内接正多边形面积 $P_n$（其实 $P_n$ 一般已不是有理数）序列的极限，而是“序列”本身，这一些读者们应该都懂得了。但是我们的直觉还能不能保存？说算术化就是完全抛弃几何直观，至少在心理上是做不到的。例如说，$\lim P_n=\pi r^2$ 至少在心理上意味着内接正多边形边数越多，$P_n$ 越接近圆的面积，这一点是无法排除的。因此，说实数 $a$ 由有理数的柯西序列 $\{a_n\}$ 决定，在我们的直观上仍旧是 $\lim a_n=a$。算术化以后能否保持这个直觉的理解呢？首先要在实数系 $\mathbb R$ 中把 $\lim a_n=a$ 讲清楚，这就是任给一个实数 $\varepsilon>0$（不再是有理数 $\varepsilon>0$ 了），必可找到正整数 $N$，使当 $n>N$ 时，$a_n$ 与 $a$ 之距离小于 $\varepsilon$。可是 $a_n$ 是有理数，$a$ 是实数，怎样讨论它们的距离呢？这就需对有理数有另一个理解：有理数就是常驻序列 $\{a_n,a_n,\ldots,a_n,\ldots\}$。于是我们有一个一般的柯西序列 $a=\{a_1,a_2,\ldots,a_n,\ldots\}$，一个常驻序列 $\{a_n,a_n,\ldots,a_n,\ldots\}$，求它们的差，再与 $\varepsilon$（为简单计，设 $\varepsilon$ 是有理数，所以是一个常驻序列 $\{\varepsilon,\varepsilon,\ldots\}$）比较大小。这样就会有下面的结果。

\noindent\textbf{基本定理 I$'$}\quad
任一个有理数柯西序列 $\{a_1,\ldots,a_n,\ldots\}$ 必收敛于它所定义的实数。

\noindent\textbf{证}\quad
$a_n-a=\{a_n-a_1,a_n-a_2,\ldots,a_n-a_m,\ldots\}$。这里 $n$ 是固定的，而 $m$ 则取一切正整数值。由于 $\{a_1,a_2,\ldots,a_m,\ldots\}$ 是柯西序列，对任意 $\varepsilon$ 必可找到 $N(\varepsilon)$，使当 $m,n>N(\varepsilon)$ 时，$|a_n-a_m|<\varepsilon/2$。现在在 $a_n-a$ 的序列表示中，$n>N(\varepsilon)$ 已是固定的，而当 $m>N(\varepsilon)$ 时，又有
\[
 -\frac\varepsilon2<a_n-a_m<\frac\varepsilon2.
\]
所以
\[
 \varepsilon-(a_n-a_m)>\varepsilon-\frac\varepsilon2=\frac\varepsilon2,
\]
\[
 \varepsilon+(a_n-a_m)>\varepsilon-\frac\varepsilon2=\frac\varepsilon2.
\]
作为定义实数的柯西序列来看，就是
\[
 \{\varepsilon,\ldots,\varepsilon,\ldots\}-(a_n-a)>0,
 \qquad
 \{\varepsilon,\ldots,\varepsilon,\ldots\}+(a_n-a)>0.
\]
亦即
\[
 \varepsilon>a_n-a,\qquad -\varepsilon<a_n-a.
\]
或者说，当 $n>N(\varepsilon)$ 时
\[
 |a_n-a|<\varepsilon.
\]
定理证毕。

\noindent\textbf{推论1}\quad
有理数系 $\mathbb Q$ 在实数系中处处稠密。

\noindent\textbf{证}\quad
所谓处处稠密就是指：任给一个实数 $a$，必可找到一个有理数列 $\{a_n\}$，使 $\lim a_n=a$。但上面证明的就是这件事。

这里附带提醒一下，当我们说“找到一个有理数序列 $\{a_n\}$”时，这里的 $a_n$ 可以按朴素的有理分数 $q/p$ 的意义来理解，而说 $\lim a_n=a$ 时，$a_n$ 一定是指常驻序列 $\{a_n,a_n,\ldots,a_n,\ldots\}$ 的等价类。

这个推论的重要性在于：它告诉我们，$\mathbb Q$ 中“空隙”太多了，比所有有理数还多，因为有理数只是可数无穷多个；尽管如此，仍然处处有有理数。现在的问题是用以上方式是否又会产生新的空隙？因此是否仍有必要再令实数组成的某些柯西序列为另一类新数，把 $\mathbb R$ 再扩充一些呢？下面的基本定理表明不再会有这样的情况：

\noindent\textbf{基本定理 I}\quad
任一个实数组成的柯西序列 $\{a_n\}$ 必定义一个实数 $a$，使得 $\lim a_n=a$。

\noindent\textbf{证}\quad
由推论1，对 $\{a_n\}$ 中的每一个实数 $a_n$，都可以找到一个有理数 $a'_n$，使 $|a_n-a'_n|<1/n$，这样一来，
\[
 |a'_n-a'_m|\le |a'_n-a_n|+|a_n-a_m|+|a_m-a'_m|
 <\frac1n+\frac1m+|a_n-a_m|.
\]
对任一 $\varepsilon>0$，由 $\{a_n\}$ 为柯西序列可知，一定有 $N(\varepsilon)$ 在，当 $n,m>N(\varepsilon)$ 时 $|a_n-a_m|<\varepsilon/3$，我们还可以取 $N(\varepsilon)>3/\varepsilon$，于是当 $n,m>N(\varepsilon)$ 时还有 $1/n,1/m<\varepsilon/3$，因此当 $n,m>N(\varepsilon)$ 时
\[
 |a'_n-a'_m|<\frac1n+\frac1m+|a_n-a_m|<\varepsilon.
\]
因此 $\{a'_n\}$ 是一个有理数的柯西序列，由基本定理 I$'$ 及推论1即有
\[
 \lim_{n\to\infty}(a_n-a)
 =\lim_{n\to\infty}(a_n-a'_n)+\lim_{n\to\infty}(a'_n-a)=0.
\]
基本定理 I 告诉我们，如果把没有极限的柯西序列称为“空隙”，则经过以上的完备化手续以后，实数系 $\mathbb R$ 中再也没有“空隙”了，因此是完备的，而且有理数系 $\mathbb Q\subset\mathbb R$ 在 $\mathbb R$ 中稠密。有没有在其它意义下产生的空隙呢？以下我们会证明，在一定意义下再也不会有空隙了，因此 $\mathbb R$ 确实是完备的。

\noindent\textbf{2. 完备性和实数系的一些基本定理}\quad
微积分中有一些基本定理，其实都是想证明有一个具某种性质的实数存在。如果证明不了，我们就会想，能不能再用上面的方法来定义新数来扩充 $\mathbb R$ 呢？如果这样，就可以说 $\mathbb R$ 中仍有空隙。这些定理的证法，最后都是归结为寻求一个柯西序列，并利用基本定理 I 得知有一个实数 $a$ 存在，然后再证明 $a$ 就是所寻求的有某种性质的实数。这样的定理我们可以举出三个来，并分别利用基本定理 I 加以证明。

\noindent\textbf{定理2（区间套定理）}\quad
设有实数闭区间套，即闭区间之序列 $\{[a_n,b_n]\}$，而且适合
\[
 [a_n,b_n]\subset[a_{n-1},b_{n-1}],
\]
若这些区间之长 $(b_n-a_n)\to0$，则必存在唯一一实数 $c$，属于所有区间 $[a_n,b_n]$ 之中。

\noindent\textbf{证}\quad
任给一个 $\varepsilon>0$，必有正整数 $N(\varepsilon)$ 存在，使当 $n>N(\varepsilon)$ 时，$b_n-a_n<\varepsilon$。现在固定这样一个 $m>N(\varepsilon)$，则 $a_{m+1},a_{m+2},\ldots$ 均含于 $[a_m,b_m]$ 中，因此
\[
 |a_{m+k}-a_m|\le |b_m-a_m|<\varepsilon,
\]
这样 $\{a_n\}$ 形成一个实数的柯西序列，因而由基本定理 I，$\lim a_n=c$ 存在。很容易证明 $a_m\le c\le b_m$，即 $c\in[a_m,b_m]$。而由于 $m$ 可以任意选取，只要 $m>N(\varepsilon)$ 即可，选 $m+k$ 代替 $m$ 易见 $c\in[a_{m+k},b_{m+k}],k=1,2,\ldots$。又由于
\[
 [a_1,b_1]\supset\cdots\supset[a_{m-1},b_{m-1}]\supset[a_m,b_m],
\]
所以 $c$ 也在 $[a_1,b_1],\ldots,[a_{m-1},b_{m-1}]$ 中，总之 $c$ 属于一切 $[a_n,b_n]$ 中。这样的 $c$ 必是唯一的，因为若有另一个 $c'\in[a_n,b_n],n=1,2,\ldots$，则因 $c'\ne c$，$|c'-c|$ 为一确定的正数，而且
\[
 |c'-c|\le |b_n-a_n|
\]
对一切 $n$ 均成立，但 $b_n-a_n\to0$，这与 $c'\ne c$ 矛盾。

\noindent\textbf{定理3（单调有界序列收敛定理）}\quad
设 $\{a_n\}$ 是单调上升的有界序列：
\[
 a_1\le a_2\le\cdots\le M.
\]
则必有唯一实数 $c$ 是 $\{a_n\}$ 的极限，$\lim a_n=c$。

\noindent\textbf{证}\quad
现在我们用反证法证明 $\{a_n\}$ 是柯西序列。若不然，对于某一个 $\varepsilon_0>0$，不论 $N$ 取多么大，其后必可找到 $a_n$ 和 $a_{n'}$（$n'>n$）使
\[
 a_{n'}-a_n=|a_{n'}-a_n|>\varepsilon_0.
\]
按此法任取一个 $N_1$，于是有 $a_{n_1}$ 和 $a_{n'_1}$（$n'_1>n_1$）使 $a_{n'_1}-a_{n_1}>\varepsilon_0$。然后取 $N(\varepsilon)=N_2>n'_1$，使其后有 $a_{n_2}$ 和 $a_{n'_2}$ 使 $a_{n'_2}-a_{n_2}>\varepsilon_0$。这样必可找出 $a_{n_k}$ 和 $a_{n'_k}$（$n'_k>n_k$）而 $a_{n'_k}-a_{n_k}>\varepsilon_0$，而且
\begin{equation}
 n_1<n'_1<n_2<n'_2<\cdots<n_k<n'_k<\cdots.
 \tag{2}
\end{equation}
于是
\[
 \begin{aligned}
 a_{n'_k}={}&(a_{n'_k}-a_{n_k})+(a_{n_k}-a_{n'_{k-1}})+\cdots\\
 &+(a_{n'_2}-a_{n_2})+(a_{n_2}-a_{n'_1})+(a_{n'_1}-a_{n_1})+a_{n_1}\\
 \ge{}&(a_{n'_k}-a_{n_k})+(a_{n'_{k-1}}-a_{n_{k-1}})+\cdots+(a_{n'_1}-a_{n_1})+a_{n_1}.
 \end{aligned}
\]
这是因为由(2)，例如 $a_{n_k}\ge a_{n'_{k-1}}$ 等等，所以上式中每一个括号均非负数。由于已设 $a_{n'_j}-a_{n_j}>\varepsilon_0$，所以由上式
\[
 a_{n'_k}\ge k\varepsilon_0+a_{n_1}.
\]
令 $k\to+\infty$，即有 $a_{n'_k}\to+\infty$ 而与 $\{a_n\}$ 是有界序列矛盾。由基本定理 I 有
\[
 \lim_{n\to\infty}a_n=c.
\]
定理证毕。

当然，对于有界的单调下降序列，本定理也成立。

\noindent\textbf{定理4}\quad
有界集必有确界。

\noindent\textbf{证}\quad
我们只证有界集 $A$ 必有下确界 $\inf A=m$。为方便计，用十进小数来表示 $A$ 的下界，于是先设 $m_0$ 是整数，$m_0$ 是 $A$ 的下界，但 $m_0+1$ 不是下界；取 $m_1$ 为0到9间的某个数字，使 $m_0+m_1/10$ 是 $A$ 的下界，但是 $m_0+(m_1+1)/10$ 不是下界。确定了 $m_1$ 以后再看小数点后的下一位，再取一个整数 $0\le m_2\le9$，使
\[
 m_0+\frac{m_1}{10}+\frac{m_2}{10^2}
\]
是 $A$ 的下界，但
\[
 m_0+\frac{m_1}{10}+\frac{m_2+1}{10^2}
\]
不是下界。仿此，可以得到一个无尽十进小数
\[
 m_0.m_1m_2\cdots m_k\cdots
\]
（若从某个 $k$ 起 $m_{k+1}=m_{k+2}=\cdots=0$，我们实际上得到的是一个有尽小数，但是我们仍然视这为无尽小数）。它的任何“一段” $m_0.m_1\cdots m_k$ 都是下界，而 $m_0.m_1\cdots(m_k+1)$ 不是下界（不过如果 $m_k=9,m_k+1=10$，就表示小数中要进一位）。但是无尽十进小数必是一个柯西序列。因此由基本定理 I，它表示一个实数 $c$，今证 $c$ 是 $A$ 的下确界。

任取 $A$ 中一元 $a\in A$，由于 $m_0.m_1\cdots m_k$ 是下界，所以
\[
 m_0.m_1\cdots m_k\le a.
\]
令 $k\to+\infty$ 即有
\[
 c\le a.
\]
所以 $c$ 是 $A$ 的下界。为了证明它是下确界，我们再取任意数 $d>c$，今证 $d$ 必不是下界。因为 $d$ 的十进小数 $n_0.n_1\cdots n_k\cdots$ 中至少有一位大于 $c$ 的十进小数的相应数字，设 $n_k$ 中第一个大于 $m_k$ 的是 $n_k$，即 $n_k\ge m_k+1$，所以
\[
 d\ge n_0.n_1\cdots n_k\ge m_0.m_1\cdots(m_k+1).
\]
但是右边必非 $A$ 之下界，所以 $d$ 也不是 $A$ 的下界，于是 $c=\inf A$。

同样可以证明上有界集必有上确界存在。

还有一些重要的定理，因为它们是关于实数空间 $\mathbb R$ 的另一重要问题——紧性问题的，我们将在下面再讲。我们现在还是把涉及完备性的话说完。

定理2—4与基本定理 I 是等价的：即是说，不但可以由基本定理 I 证明定理2—4，其实由定理2—4的任何一个也可以证明基本定理 I 和另两个定理。其原因是，这些结果都是以不同方式回答了一个问题：有理数系 $\mathbb Q$ 的空隙应该怎样去刻画并填补起来。例如区间套定理讲的就是：一个有理数区间套若其各区间长度趋向0，可能会出现在有理数系 $\mathbb Q$ 中找不到一个公共元，这样就出现了空隙，而定理3告诉我们：没有关系，有了这样的区间套，就有一个柯西序列例如 $\{a_n\}$，它定义一个新数（实数）$a$，把这个 $a$ 填到空隙里去，这个区间套就有了唯一的公共点 $a$，而区间套定理就得到了证明。很自然，我们也可以走另一条路，如果某个有理数区间套出了空隙——即不能以某有理数为公共点，就说它定义了一个新数，称为实数。把这个实数“塞”进区间套的每个区间内，于是此区间套就有了一个公共点（即此实数）。这样定义了的实数系仍可进行四则运算，仍有大小（还可加上阿基米德公理），于是此区间套可以说是趋向这个实数的，而原来的有理数系按原有的大小说也好，按扩充成的实数系的大小来说也好都在此实数系中稠密（相应于基本定理 I$'$ 和推论1）。如果再用扩充成的实数系作区间套，就不会再出现空隙，而如它必含有唯一公共点（相应于基本定理 I，即定理2）。把它称为基本定理 III，则由它不但可以证明定理3、4，还可以证明关于极限存在的柯西准则：任一实数柯西序列必收敛于唯一实数。总之这几个定理都可以作为基本定理，但它的证明必须分成两步：先是基本定理 X$'$：这是由有理数系 $\mathbb Q$ 填补空隙而得实数系 $\mathbb R$；然后再是基本定理 X，这是说实数系 $\mathbb R$ 中再没有空隙了。这就是说，上面的基本定理和三个定理实际上给出了构造实数系 $\mathbb R$ 的等价的方法。

于是还有一个问题：会不会在某种新环境下 $\mathbb R$ 出现新的空隙，需要用新的方法填补，而得到一个与上面作出的 $\mathbb R$ 不同的实数系 $\widetilde{\mathbb R}$ 呢？例如在通常较详细的微积分教本上都是用的戴德金分割来定义实数系 $\mathbb R$ 而与以上诸定理均不同。我们甚至应想到，还会不会有其它方法作出 $\mathbb Q$ 的不同扩充来呢？下面的定理告诉我们：不会。

\noindent\textbf{基本定理 II}\quad
有理数系 $\mathbb Q$ 扩充为完备的实数系并使 $\mathbb Q$ 在其中稠密的方法在等距同构意义下是唯一的。

\noindent\textbf{证}\quad
先把定理的含意解释清楚。设我们已把有理数系 $\mathbb Q$ 扩充为实数系 $\widetilde{\mathbb R}$，而且使 $\widetilde{\mathbb R}$ 中的运算关系与次序（即大小）关系与 $\mathbb Q$ 中原有的运算关系、次序关系相协调；例如设 $a\in\mathbb Q\subset\widetilde{\mathbb R}$，则 $a$ 在 $\mathbb Q$ 中的大小（即原来的大小）与它在 $\widetilde{\mathbb R}$ 中的大小一致。$\mathbb Q$ 在 $\widetilde{\mathbb R}$ 中稠密，即 $\widetilde{\mathbb R}$ 中任一元 $r$ 必可用 $\mathbb Q$ 中的元之序列 $\{a_n\}$ 去逼近：$|a_n-r|\to0$。这里的距离既可以说是 $\mathbb Q$ 中的距离，也可以说是 $\widetilde{\mathbb R}$ 中的距离。这个定理说，如果有两种适合这样的要求的扩充 $\widetilde{\mathbb R}_1$ 与 $\widetilde{\mathbb R}_2$，则二者必等距同构，即是说 $\widetilde{\mathbb R}_1$ 中的“实数”与 $\widetilde{\mathbb R}_2$ 中的“实数”，无论就其大小而言或者就其运算关系而言都是一样的。但是实数除了是运算对象与大小关系而外还有什么含义呢？所以我们说 $\widetilde{\mathbb R}_1$ 和 $\widetilde{\mathbb R}_2$ 是一样的。因此，$\mathbb Q$ 的扩充是唯一的。明白了定理的含意后，证明是很容易的。我们记上面讲的用柯西序列定义作扩充的实数系为 $\mathbb R$。今证任何适合定理要求的扩充 $\widetilde{\mathbb R}$ 必与 $\mathbb R$ 等距同构。

设有 $\tilde r\in\widetilde{\mathbb R}$，则由 $\mathbb Q$ 在 $\widetilde{\mathbb R}$ 中稠密，必可找到 $\mathbb Q$ 中一个序列 $\{a_n\}$，$a_n\to\tilde r$，这个 $\{a_n\}$ 一定是柯西序列。因此由 $\mathbb R$ 之定义，$\{a_n\}$ 定义一个元 $r\in\mathbb R$，我们令 $\tilde r$ 对应于 $r$。反过来也一样，任给一个 $r\in\mathbb R$ 也必有一个 $\tilde r\in\widetilde{\mathbb R}$ 与它对应。所以 $\widetilde{\mathbb R}$ 与 $\mathbb R$ 之间是一一对应的。而且若 $r\sim\tilde r,s\sim\tilde s$，则 $r+s\sim\tilde r+\tilde s$，$rs\sim\tilde r\tilde s$，而且域中的运算在此对应下是不变的。所以 $\widetilde{\mathbb R}$ 与 $\mathbb R$ 是同构的。$r>s$ 必有 $\tilde r>\tilde s$，所以次序也不变。$\widetilde{\mathbb R}$ 中的距离关系与 $\mathbb R$ 中的距离关系也是相同的，即 $\widetilde{\mathbb R}$ 与 $\mathbb R$ 等距。证毕。

至此，我们明确地证明了，只有一个实数系 $\mathbb R$，它是一个适合阿基米德公理的、完备的、有序的（代数）域，而有理数域 $\mathbb Q$ 是它的处处稠密的子集。上面全部讨论集中在其完备性上。

\noindent\textbf{3. 紧性}\quad
紧性的讨论始自19世纪末、20世纪初，比完备性的讨论晚了许多。最早看出它的重要性的数学家大概是博雷尔（E. Borel），他是从这样一个事实开始的。

\noindent\textbf{定理5（海涅—博雷尔（Heine--Borel）定理）}\quad
设 $[a,b]$ 是 $\mathbb R$ 中的一个闭区间，$\{U_\lambda\}$ 是它的一个开覆盖，则从 $\{U_\lambda\}$ 中一定可以取出有限多个 $U_1,\ldots,U_k$，构成 $[a,b]$ 的一个子覆盖：
\begin{equation}
 [a,b]\subset\bigcup_{i=1}^k U_i.
 \tag{3}
\end{equation}

\noindent\textbf{证}\quad
在证明此定理之前，我们还是先来解释一下它的意义。“覆盖”的含意很简单，这里就是说
\[
 [a,b]\subset\bigcup_\lambda U_\lambda.
\]
需要注意的是，$\{U_\lambda\}$ 不一定是可数多个 $U_\lambda$（若 $U_\lambda$ 为可数多个叫可数覆盖，这时我们常取 $\lambda=1,2,\ldots$），所以 $\lambda$ 就只是一个参量，它的集合 $\Lambda$ 可以是有限集，也可以是可数集或不可数集。但一定不是空集，因为那样一来，就没有 $U_\lambda$ 了，也就不会覆盖什么。开覆盖即指 $\{U_\lambda\}$ 之每个元素都是开集。但我们知道，在 $\mathbb R$ 中每个开集都是最多可数多个开区间之并，所以可把组成各个 $U_\lambda$ 的开区间分开来重新编号，并仍记作 $\{U_\lambda\}$，则每个 $U_\lambda$ 又都可以设为开区间。不过，因为这个定理太重要了，所涉及的概念不只会被应用在 $\mathbb R$ 上，那时也没有开区间的概念，所以我们现在还是说 $U_\lambda$ 就是一个开集。

但是我们的证明中是还要利用到每个 $U_\lambda$ 均可分为若干个开区间。其证明其实很容易。若 $[a,b]$ 只需有限多个 $U_\lambda$ 即可覆盖，则没有什么好证明的了。若不然，将它平分为两个闭区间：
\[
 [a,b]=\left[a,\frac{a+b}{2}\right]\cup\left[\frac{a+b}{2},b\right]
\]
（两个小闭区间有一个公共点 $(a+b)/2$ 并无关系），则至少有一个需无限多个 $U_\lambda$ 才能覆盖。设这一个是 $[a_1,b_1]$，再把 $[a_1,b_1]$ 平分为二，又至少有一个需无限多个 $U_\lambda$ 才能覆盖。仿此，我们会得到一个区间套 $[a_k,b_k]$，其长为 $2^{-k}(b-a)\to0$。由区间套定理，必定有一点 $c$ 在所有 $[a_k,b_k]$ 中。因为 $\{U_\lambda\}$ 覆盖了 $[a,b]$，所以 $c$ 必在某个 $U_{\lambda_0}$ 之内，从而必在某一开区间（不妨就设 $U_{\lambda_0}$ 是一个开区间）内。由开区间的性质，必有以 $c$ 为中点的小开区间 $(c-r,c+r)\subset U_{\lambda_0}$，其长为 $2r$。但 $c$ 又在一切 $[a_k,b_k]$ 中，若某个 $[a_k,b_k]$ 之长小于 $r$，必有 $[a_k,b_k]\subset U_{\lambda_0}$ 中。只要 $k$ 充分大，$[a_k,b_k]$ 之长自然小于 $r$，所以，上述平分只要进行 $k$ 步，则根本无需无限多个 $U_\lambda$ 才能覆盖——一个 $U_{\lambda_0}$ 就够了。这就是矛盾，定理证毕。

这个定理的证明本质地利用了实数系的完备性——它用了区间套定理。这样，使人或者有这样的想法：定理5是否和前面的定理2—4一样，只不过是完备性的另一种说法呢？否，这里讲的是紧性。什么是紧性呢？要定义它需要有开集、闭集等等概念。一个集合，赋以开集概念以后，称为一个空间（我们下面要详细讲，请看第三节），一个空间 $K$ 为紧就是说，它的任意开覆盖都有一个有限的子覆盖。所以定理5讲的是：$\mathbb R$ 中的任一闭区间 $[a,b]$ 都是紧的。（许多书上更小心一点，强调指出是有界闭区间；我们没有这么小心，因为到现在为止，我们都认为无穷大不是实数——阿基米德公理告诉我们这一点，因此写成 $[a,b]$ 就意味着 $a,b$ 都是有限数而不是 $\pm\infty$，所以 $[a,b]$ 自然是有界闭区间。）因为定理讲的是关于某些实数集的性质，证明中自然会用到实数集的性质，包括完备性。这不意味着紧性这个概念是完备性概念的衍生物。

我们要注意，定理5只说 $[a,b]$ 为紧，而没有说 $\mathbb R$ 为紧。事实上，$\mathbb R$ 恰好不紧。我们可以用一个例来证明这件事，作一个开覆盖 $\{U_k\}$，这里 $k=0,\pm1,\pm2,\ldots$ 而
\[
 U_k=\left(k-\frac23,k+\frac23\right),
\]
$\mathbb R\subset\bigcup_kU_k$ 是明显的事，但是 $\{U_k\}$ 中决没有 $\mathbb R$ 的有限子覆盖。因为若要覆盖 $\mathbb R$，就需要覆盖其上的无穷多个坐标为整数的点。但每一个区间 $U_k$ 只可能含一个“整点” $k$，因此要覆盖 $\mathbb R$ 就必须用完全部的 $U_k$，而不可能有有限子覆盖。$\mathbb R$ 整个说来不紧，但每一点都有一个邻域（不妨设为一区间 $(a,b)$），其闭包 $[a,b]$ 为紧。这种情况称为局部紧。所以 $\mathbb R$ 不紧而仅只为局部紧。紧与局部紧是有重大区别的，这种区别在许多数学问题上有重大影响。

$\mathbb R$ 中与紧性有关还有几个概念，一是魏尔斯特拉斯—波尔查诺性质，一个空间 $K$ 有此性质是指此空间中任一序列 $\{x_n\}\subset K$，必有一个位于 $K$ 内的极限点，即这样的点 $x_0$：其任一邻域 $I=(x_0-\varepsilon,x_0+\varepsilon)$ 与 $K$ 之交 $I\cap K$ 中均有此序列中的点。另一是列紧性：我们说空间 $K$ 是列紧的即指 $K$ 中之任一序列 $\{x_n\}\subset K$ 必有一个收敛子序列，其极限 $x_0\in K$。这里我们有：

\noindent\textbf{定理6}\quad
设 $K\subset\mathbb R$，其中之开、闭集即指 $\mathbb R$ 中的一个开、闭集与 $K$ 之交，这时以下三个命题等价：

\noindent (i) $K$ 为紧；

\noindent (ii) $K$ 具有魏尔斯特拉斯—波尔查诺性质；

\noindent (iii) $K$ 具有列紧性。

这个定理的证明并不太难，而主要涉及一些技术性的细节，所以我们略去它。但重要的是要指出紧与列紧、与魏尔斯特拉斯—波尔查诺性质是不相同的概念，只不过对 $\mathbb R$ 的子集 $K$ 三者一致，而在更广泛的情况下就不一定一致了。若 $K=[a,b]$，则(ii)与(iii)都是微积分教程中常见的（但是重要的）结果。正因为如此，有些读者会不注意把它们区别开来。

微积分教本中通常都要讲到闭区间 $[a,b]$ 上的连续函数的一些基本性质。我们已经知道 $[a,b]$ 是 $\mathbb R$ 的一个紧集，而通常微积分教本中通常回避使用这个词，我们则坚持使用紧集这个词，除非是必须用到 $[a,b]$ 作为一个区间才有的性质时，说明我们讨论的是区间。例如下面的(iii)，其目的向读者“灌输”紧性之重要。我们将要证明下面的定理。

\noindent\textbf{定理7}\quad
设 $K\subset\mathbb R$ 是一个紧集，$f(x)$ 是 $K$ 上的连续函数，则有

\noindent (i) $f(x)$ 在 $K$ 上有界，而且必能达到其上、下确界使之成为最大、最小值；

\noindent (ii) $f(x)$ 必在 $K$ 上一致连续；

\noindent (iii) 若 $K=[a,b]$，则 $f(x)$ 可以取到其最大、最小值之间的一切值。

\noindent\textbf{证}\quad
在证明之前，先对定理作一些解释。我们上面特别提醒，紧集 $K$ 不一定是一个区间，但讲函数 $f(x)$ 在一点 $x_0$ 连续时，$x_0$ 总是取为某一开区间的内点。现在所谓 $f(x)$ 在 $x_0\in K$ 连续，是指在
\[
 \lim_{x\to x_0}f(x)=f(x_0)
\]
中 $x$ 必须在 $K$ 内趋向 $x_0$。这在通常的微积分教本中本已见过，例如说 $f(x)$ 在 $[a,b]$ 的左端 $x_0=a$ 处连续，就要求 $x$ 从 $a$ 点右方（即 $[a,b]$ 内部）趋向 $a$ 点，并称为右连续，其实就是 $f(x)$ 相对 $[a,b]$ 的连续性的意思。甚至 $K$ 还可以是孤立点（有限多个点所成的集合一定是紧的）$x_0$，这时 $\lim_{x\to x_0}f(x)=f(x_0)$ 中 $x\to x_0$ 只能解释为 $x=x_0$，所以 $f(x)$ 只要在 $x_0$ 处有定义就是 $f(x)$ 在 $x_0$ 连续。（同理，$\lim_{x\to x_0}f(x)=A$ 就可能没有意义了，因为这里既要求 $x\in K$，又要求 $x\ne x_0$，二者可能矛盾。）与此类似，当我们说到 $K$ 的开覆盖 $\{U_\lambda\}$ 时，所谓 $U_\lambda$ 是开集，是指 $U_\lambda$ 是 $\mathbb R$ 中一个开集（它是若干个开区间之并）$\widetilde U_\lambda$ 与 $K$ 之交：$U_\lambda=\widetilde U_\lambda\cap K$。甚至为简单起见我们就说 $U_\lambda$ 是 $(\alpha,\beta)$。这种问题在本书中已提到过好几次。在第三节讲到拓扑空间的子空间时我们再系统地说明。下面开始证明。

(i) 因为 $f(x)$ 是连续的，故对任一点 $x_0\in K$，对某个固定的 $\varepsilon_0$ 必有 $\delta>0$ 存在使在
\[
 U_{x_0}=(x_0-\delta,x_0+\delta)
\]
内（准确些说是在 $U_{x_0}=(x_0-\delta,x_0+\delta)\cap K$ 内），有
\begin{equation}
 f(x_0)-\varepsilon_0<f(x)<f(x_0)+\varepsilon_0.
 \tag{4}
\end{equation}
于是 $\{U_{x_0}\}$（参数 $x_0$ 现在就是 $x_0$，$\Lambda$ 就是 $K$）成为 $K$ 的一个开覆盖，而因 $K$ 为紧，可以找到一个有限子覆盖 $\{U_i\},i=1,2,\ldots,k$。在 $U_i$ 中，
\begin{equation}
 m_i=f(x_i)-\varepsilon_0<f(x)<f(x_i)+\varepsilon_0=M_i,
 \qquad i=1,2,\ldots,k.
 \tag{5}
\end{equation}
令 $M=\max_{1\le i\le k}M_i$，$m=\min_{1\le i\le k}m_i$，则知在 $K$ 上
\[
 m\le f(x)\le M.
\]
因此 $f(x)$ 有界。

既有界必有上下确界，令 $\mu$ 为其上确界，我们用反证法证明一定有一点 $x^*\in K$ 使 $f(x^*)=\mu$。如果这样的 $x^*$ 不存在，则在 $K$ 上必有 $\mu-f(x)>0$。令
\[
 \varphi(x)=\frac1{\mu-f(x)},
\]
则 $\varphi(x)$ 在 $K$ 上连续，而如上所证，$\varphi(x)$ 在 $K$ 上必有上界 $\lambda$，
\[
 \frac1{\mu-f(x)}\le\lambda.
\]
因此 $f(x)\le\mu-1/\lambda<\mu$ 在 $K$ 上成立，因此 $\mu$ 不是上确界，命题得证。

我们要看一下紧性在这个证明中起的作用，关键在于(4)式本来对一切 $U(x_0)$ 成立，所以在任一 $U_{x_0}$ 中有 $m(x_0)\le f(x)\le M(x_0)$。但只要有无穷多个 $x_0$ 点，就有可能有无穷多个不同的 $M(x_0)$，由其中不一定能找出一个最大的（有时为方便起见我们也说其最大是 $+\infty$），紧性则把无穷多个(4)归结为有限多个(5)，而由有限多个 $M_i$ 中自然有最大的 $M$（而自然是指有限的实数）存在。所以紧性的作用在于能把无穷多个对象归结为有限多个对象。人们常说，紧集是“最接近于有限集”的无穷集合就是这个意思。

上面说有时为方便起见常说无上界的函数就是以 $+\infty$ 为上界的函数，而且这时上确界就是 $+\infty$。按阿基米德公理，不可能有一个实数是无穷大，所以这只是一个方便的说法，但我们时常要用这种说法，本节之末我们再稍作解释。

(ii) 对任意 $\varepsilon>0$，对 $K$ 中每一点 $x^*$ 必有 $\delta'(\varepsilon,x^*)$ 存在，使当
\[
 x\in(x^*-\delta',x^*+\delta')\cap K
\]
时，
\[
 |f(x)-f(x^*)|<\frac\varepsilon2.
\]
在此区间中任取两点 $\bar x,x$，则有
\begin{equation}
 |f(\bar x)-f(x)|\le|f(\bar x)-f(x^*)|+|f(x)-f(x^*)|<\varepsilon.
 \tag{6}
\end{equation}
把区间 $(x^*-\delta',x^*+\delta')$ 缩小一半，并记
\[
 U_{x^*}=\left(x^*-\frac{\delta'}2,x^*+\frac{\delta'}2\right)\cap K,
 \qquad x^*\in K,
\]
则 $\{U_{x^*}\}$ 仍是 $K$ 的开覆盖，所以可以从其中选出有限多个：
\[
 U_i=\left(x_i-\frac{\delta_i}2,x_i+\frac{\delta_i}2\right)\cap K,
 \qquad i=1,2,\ldots,k,
\]
仍是 $K$ 的开覆盖。取
\[
 \delta=\min_{1\le i\le k}\frac{\delta_i}{2}>0.
\]
并且证明，$K$ 中任意两点 $x,\bar x$，只要 $|x-\bar x|<\delta$，则这两点一定位于同一个 $(x_i-\delta_i,x_i+\delta_i)$ 中。事实上，既然 $\{U_i\}$ 是 $K$ 的覆盖，则 $\bar x$ 必在某一个 $U_i$ 中，不妨设在 $U_1$ 中。因此
\[
 |\bar x-x_1|<\delta_1/2,
\]
而
\[
 |x-x_1|\le|x-\bar x|+|\bar x-x_1|<\delta+\delta_1/2\le\delta_1,
\]
$x$ 和 $\bar x$ 同在 $(x_1-\delta_1,x_1+\delta_1)$ 中，而由(6)式
\[
 |f(x)-f(\bar x)|\le|f(x)-f(x_1)|+|f(x_1)-f(\bar x)|<\varepsilon.
\]
所以 $f(x)$ 在 $K$ 上一致连续。

(iii) 的证明全不相同。它实际上与紧性无关，但我们还是先利用紧性来证明它。首先我们已知 $f(x)$ 必在两点达到最大值 $M$ 与最小值 $m$：$f(a)=M,f(b)=m$。我们不妨设 $a<b$，并且用 $[a,b]$ 代替原来的区间。我们要证对任一 $\mu$，$M>\mu>m$，在 $[a,b]$ 中有一点 $c$ 使 $f(c)=\mu$。为此，用 $\varphi(x)=f(x)-\mu$ 代替 $f(x)$，则有 $\varphi(a)>0,\varphi(b)<0$，因此我们把待证明的命题改为：若 $f(x)$ 在 $[a,b]$ 的两个端点异号，则一定存在一点 $c\in[a,b]$，使 $f(c)=0$。我们仍用反证法，设这样的 $c$ 不存在。任取一点 $x^*\in[a,b]$，必有 $f(x^*)\ne0$。因为 $f(x)$ 是连续的，故当 $\delta$ 充分小时，$f(x)$ 在 $U_{x^*}=(x^*-\delta,x^*+\delta)$ 中与 $f(x^*)$ 符号相同。$\{U_{x^*}\},x^*\in[a,b]$ 构成 $[a,b]$ 的一个开覆盖，于是它有一个有限的子覆盖 $U_1,U_2,\ldots,U_k$。我们不妨将它从左到右排列，而且假设它们都是开区间，这样 $a\in U_1,b\in U_k$，而且相邻的两个 $U_i$ 必定相交：$U_{i-1}\cap U_i=I_i$ 是一个开区间。在每个 $U_i$ 上，$f(x)$ 都是不变号的，于是在 $U_1$ 中 $f(x)$ 与 $f(a)$ 同号。在 $I_2$ 中取一点 $x_2$，则 $x_2\in U_1$，从而 $f(x_2)$ 与 $f(a)$ 同号，但是在 $U_2$ 中 $f(x)$ 也是不变号的，而且 $x_2\in U_2$，所以在 $U_2$ 中 $f(x)$ 也与 $f(a)$ 同号。仿此，在所有的 $U_1,\ldots,U_k$ 中 $f(x)$ 均与 $f(a)$ 同号，从而 $f(b)$ 与 $f(a)$ 同号，而这与假设矛盾，故命题得证。

这个命题与前两个有何不同？我们仔细分析一下。

首先，它利用了紧性，这样我们只要推理有限步即可得到结论。这与前面说的紧集是最接近有限集的无限集合是一样的。其次证明中利用 $f(x)$ 连续而且 $f(x)\ne0$，因此，只要 $\delta$ 充分小，$f(x)$ 在 $U_i$ 中不变号。造成矛盾的不只是 $f(x)$ 不变号，更重要的是，$U_1,\ldots,U_k$ 两两相连，有公共点 $x_i\in U_{i-1}\cap U_i$，这样才会造成矛盾。若不是讨论 $[a,b]$，而是一般的 $K$，这种两两相连就可能不出现，例如设 $K=[0,1]\cup[2,3]$，它是一个紧集，于是由 $\{U_i\}$ 中可以找到有限子覆盖 $U_1,\ldots,U_k$，很可能 $U_1\cup U_2\cup U_3$ 覆盖了 $[0,1]$，$U_4\cup\cdots\cup U_k$ 覆盖了 $[2,3]$，而 $U_3$ 与 $U_4$ 可能不相交。所以我们利用了 $[a,b]$ 的另一个重要特征——连通性——连通性与完备性、紧性一样是非常重要的性质，许多重要的数学结果（例如这个命题）来自连通性。因篇幅的原因，我们这里不再多说，§3将专门讨论有关连通性问题。

总之关于实数系 $\mathbb R$，我们现在知道了它是一个有序的域，它满足阿基米德公理，而且它具有完备性、局部紧性与连通性。

最后我们讲一下关于 $\pm\infty$ 的问题。阿基米德公理已经告诉我们，它们不是实数，但是为了方便起见，有时我们还是把它们“当数看待”，并作如下规定：
\[
 (\pm\infty)+(\pm\infty)=\pm\infty;
\]
\[
 a+(\pm\infty)=\pm\infty;
\]
\[
 a\cdot(\pm\infty)=\pm\infty\quad\text{若 }a>0;
\]
\[
 a\cdot(\pm\infty)=\mp\infty\quad\text{若 }a<0;
\]
\[
 (\pm\infty)\cdot(\pm\infty)=+\infty,\qquad (\pm\infty)\cdot(\mp\infty)=-\infty.
\]
$(\pm\infty)-(\pm\infty)$ 无定义，$0\cdot(\pm\infty)$ 有的书上说是无定义，有的书上则规定 $0\cdot(\pm\infty)=0$。

引进 $\pm\infty$，并且得到所谓“扩充的实数系”，纯粹只是为了讲话方便，例如上无界集就说是以 $+\infty$ 为上界，它的上确界也是 $+\infty$。这种语言在涉及集合、测度、积分时用得多一些，本书偶尔也会用一点。
